\scanpage{219}
\mspage{218}
$X$ \deleted{schéma} préschéma de type fini \inserted{normal} irréductible sur corps $k$

$v,w$ deux places archimédiennes de \deleted{\unsure{le}} $X$, d'où
foncteurs fibres $\xi_v^{\mathbb Z},\xi_w^{\mathbb Z}$. Question : a) \inserted{Quand} sont-ils
isomorphes ? b) \deleted{Est-il vrai que} Quand est-il
vrai que pour tte sous-catégorie \inserted{mult.} de t.f. de
$\mathbb M(X)$, les restrictions de ces foncteurs à
$\mathbb M_\alpha(X)$ sont isomorphes ?

Nous pouvons supposer, quitte à remplacer $k$ par
\deleted{\unsure{Soit déjà les}}
sa clôture alg. dans $k(X)$, que $k$ est alg.
clos dans $k(X)$ [utilise le fait $X$ \uline{normal} !],
donc $X$ est géom\textsuperscript{t} irred. sur $k$. Soit \deleted{$k_0$}
\inserted{$k_0$ la clôture alg. de $\mathbb Q$ dans $k$.}
\deleted{le corps \unsure{de définition} des places (= \unsure{sur $k_0$}).
\uline{induites sur $k_0$}
\unsure{archimédiennes}
\uline{Contient la clôture alg. $E$ de $\mathbb Q$ dans $k$} (p.ex. $k_0$ alg. fermé dans $k$).
1) $k_0$ est algébriquement fermé dans $k$. Alors}

1°) $v,w$ coïncident sur $k_0$.
\[
\text{a) }\xi_v^{\mathbb Z}|\mathbb M_\alpha(X)\simeq\xi_w^{\mathbb Z}|\mathbb M_\alpha(X)
\quad\text{pour }\forall\alpha.
\]
En effet, $X$ et $\mathbb M_\alpha(X)$ proviennent de
$X_0$, \inserted{$\mathbb M_\alpha(X_0)$}, de type fini sur $A_0$, \inserted{$k_0$}\unsure{-alg.} de
\deleted{$X_0\longrightarrow\operatorname{Spec}A_0$ fid plat \unsure{univers\textsuperscript{t}} connexe \unsure{et} $\operatorname{Spec}A_0\longrightarrow\unsure{k_0|k}$ $^{*}$}
type fini de $k_0$ \deleted{$\unsure{E}$}. On est ramené
au cas où $k=k_0$, où on obtient le
résultat voulu grâce au groupoïde fondamental
transcendant. \deleted{\unsure{Alors}}

\margin{\diagram{\begin{tikzcd}[row sep=small,column sep=small]
&X_0\arrow[r,no head]\arrow[d,no head]&X\arrow[d,no head]\\
k_0\arrow[r,no head]&A_0\arrow[r,no head]&k
\end{tikzcd}}}

\margin{\deleted{\diagram{\begin{tikzcd}[row sep=small,column sep=small]
&&&\unsure{X_0}\arrow[r,no head]&X\arrow[dd,no head]\\
&&X_0\arrow[d,no head]&&\\
E\arrow[r,no head]\arrow[drr]&k_0\arrow[r,no head]\arrow[ur,no head]&\unsure{A_0}\arrow[rr,no head]&&k\\
&&A_0\arrow[uuur,no head]&&\\
&\mathbb Q\arrow[uu,no head]\arrow[rr,no head]&&E&
\end{tikzcd}}}}

b) Si \deleted{$\unsure{k=k_0}$}\inserted{$v$ et $w$ coïncident sur $k$}, cet argument montre que
$\xi_v^{\mathbb Z}\simeq\xi_w^{\mathbb Z}$.
\inserted{Sinon,} \unsure{mais rien ne prouve qu'on ait}
\deleted{Si $k\ne k_0$, \unsure{je pense que} $\xi_v^{\mathbb Z}\ne\xi_w^{\mathbb Z}$ ; \illegible{2}}

\margin{$^{*}$ on peut supposer $X_0\longrightarrow\operatorname{Spec}A_0$ fid plat \unsure{univers\textsuperscript{t} connexe}, et
$\operatorname{Spec}(A_0)\mathrel{\text{---}}k_0$ \unsure{géom. connexe}, donc $X_0\mathrel{\text{---}}\operatorname{Spec}k_0$ \unsure{géom. connexe}.}
